\section{极值格截面与设计矩}
\label{sec:sections}

前两章直接给出构型的每一层，再检查层间内积。本章从已知的大型格壳中选出
满足内积条件的向量，并通过投影得到低维构型。需要计算的不是整个母壳，
而是若干投影特征所对应的向量数。43维计算与一个五维截面正交的壳向量；
45维选取一个三维截面的三个纤维，再投影到正交补。两者都将几何选择转为
有限的重数问题。

\subsection{母壳、投影域与矩恒等式}

48维极值偶幺模格的最小平方范数为6。其 theta 级数为

\[
E_4^6-1440E_4^3\Delta+125280\Delta^2
=1+52416000q^3+O(q^4),
\]

所以极小壳含52416000个点。这里
\(E_4=1+240q+2160q^2+6720q^3+\cdots\)，\(\Delta=q-24q^2+252q^3+\cdots\)；权24模形式的常数项及前两个零系数确定上述表达式
\cite{conway1999sphere}。Venkov 定理保证这个壳为球面11-design
\cite[定理4.1]{nebe-designs}。

球面11-design意味着次数不超过十一的多项式在这个有限壳上的平均值，
与在相同半径球面上的平均值一致。因此，即使尚未逐一列出壳向量，
也能够写出它们沿给定方向的低次矩。例如，
\[
\sum_{x\in X}(x\cdot t)^2=\frac{6|X|}{48}\|t\|^2,\qquad
\sum_{x\in X}(x\cdot t)^4=\frac{3\cdot6^2|X|}{48\cdot50}\|t\|^4.
\]
将向量按投影分组后，等式左边成为未知重数的线性组合，右边则完全已知。

取截面基向量 \(t_i\)，Gram 矩阵为 \(H\)。壳向量的签名为
\(y_i=x\cdot t_i\)，其截面投影的平方范数为
\(y^{\mathsf T}H^{-1}y\le6\)。若截面向量属于母格，则这些签名为整数；与截面中已知短向量的点积界进一步缩小候选域。对次数
\(2m\le10\)，球面矩公式给出

\[
\sum_{x\in X}\prod_{\ell=1}^{2m}(x\cdot t_{i_\ell})
=\frac{|X|6^m}{48\cdot50\cdots(48+2m-2)}
\sum_{\mathcal P}\prod_{(a,b)\in\mathcal P}H_{i_ai_b},
\]

其中 \(\mathcal P\)
遍历这些位置的所有配对；次数零的式子为总点数。因而投影签名的未知重数满足有理线性方程。43维用这些方程精确求出目标统计量；45维再结合重数非负性，得到一个目标统计量的下界。

\subsection{43维：由母截面确定子截面的点数}

设 \(X\) 是 \(P_{48n}\) 的最小壳 \cite{latticecatalogue}。给定的整数矩阵将 \(\sqrt3E_8\)
嵌入母格，其中前五个简单根张成 \(\sqrt3D_5\) 截面 \(U\)。集合
\(X\cap U^\perp\) 位于 43 维空间。

根的编号由以下标准坐标固定：
\[
\alpha_1=\tfrac12(1,-1,-1,-1,-1,-1,-1,1),\qquad
\alpha_2=e_1+e_2,\qquad
\alpha_j=e_{j-1}-e_{j-2}\quad(3\le j\le8).
\]
所给 $8\times48$ 整数矩阵 $T$ 与母格 Gram 矩阵 $G$ 满足
$(TGT^{\mathsf T})_{ij}=3\alpha_i\cdot\alpha_j$；第1至5行指定所用子截面。
这同时固定了嵌入位置和抽象的 $D_5$ 根型。嵌入 $\sqrt3E_8$ 母截面的路线
见 Dorofeev--Sun--Wang~\cite{dorofeev}；这里利用其投影数据计算这个
$D_5$ 的正交部分。

\begin{proposition}[指定的 $D_5$ 正交部分]
\label{prop:section43}
对证书给出的 $G,T$，壳 $X\cap U^\perp$ 恰有2553792个向量，因而
$K(43)\ge2553792$。
\end{proposition}

这里无需改变选中向量的方向。不同极小向量满足
$\|x-y\|^2\ge6$，因此 $x\cdot y\le3$；除以 $\sqrt6$ 即为单位球面上的
接吻构型。问题集中为如何精确计算 $|X\cap U^\perp|$。

每个最短向量在母截面上的投影可写成 \(\lambda/\sqrt3\)，其中
\(\lambda\in E_8\)。范数约束及它与每个嵌入根的内积约束，给出完整的
26,641 个投影签名，其中 240 个是截面根本身的例外签名。

这是对全部可能投影的枚举。非例外的 $\lambda$ 满足 $\|\lambda\|^2\le18$，
且对每个 $E_8$ 根 $r$ 有 $|\lambda\cdot r|\le3$。
枚举所得平方范数为0、2、4、6、8的五层大小依次为
1、240、2160、6720、17280；再加入240个 $\lambda=3r$ 的例外，得到上述完整域。

极值 48 维偶幺模格的最小壳是
11-design，因此次数不超过十的偶次矩等于相应球面矩。对母截面的 Weyl
群平均，五个非例外范数轨道上的平均重数为

\[
1092000,\quad116235,\quad9180,\quad495,\quad63/4.
\]

这些数是纤维点数在轨道上的平均值，因此可以是有理数。

目标正交子空间在这五个轨道上分别含 \(1,12,6,24,0\)
个签名，另有十二个例外根，所以平均接触数是

\[
1092000+12(116235)+6(9180)+24(495)+12=2553792.
\]

这已经说明至少一个子截面含有不少于 2,553,792
个点。为了精确确定所给子截面的点数，再只对保持该子截面的子群平均。
具体取 $-I$、根 $\alpha_1,\ldots,\alpha_5$ 的反射，以及与这五个根均正交的
十二个 $E_8$ 根的反射作为生成元。后十二个根构成 $A_3$，反射公式均为
$s_r(x)=x-(x\cdot r)r$。这些公式直接确定所需群，无须给出母格自同构。
完整签名域在该群下分成43个轨道；有理矩方程记为 \(Av=b\)，目标计数记为 \(c^tv\)。证书提供
\(w\) 满足

\[
A^tw=c,\qquad w^tb=2553792.
\]

这就证明了命题~\ref{prop:section43}。平均所需的是签名域、矩方程与目标统计量被保留，不要求这个子群一定延拓为整个母格的自同构。

这个计算利用的是线性目标的可确定性，而不是所有纤维的可确定性。
即使 $Av=b$ 有多个解，只要目标行 $c^{\mathsf T}$ 位于 $A$ 的行空间中，
$c^{\mathsf T}v$ 就在所有解上取同一个值。研究时可先用模数线性代数筛选
独立矩行，再求出有理乘子；读者复算时只需检查上述两个有理恒等式。
完整签名域在这里仍是关键：全壳的矩右端对应全部投影，正交条件应写在
目标向量 $c$ 中，而不是先丢弃其余纤维再使用同一个右端。

比较构造中的截面具有相同的抽象 \(D_5\) Gram
型。改变具体嵌入可以改变正交壳层的点数；这一嵌入可延拓到 \(E_8\)
母截面，使额外投影约束能够认证该点数。增加方程是在确定一个固定集合的大小，
而不是改变这个集合。

\subsection{45维：用298个见证控制投影点数}

这里使用一个明确指定的 48
维二次剩余偶邻格。附录\ref{app:ambient}通过三元码的最小重量、特征向量和奇偶关系证明该格的最小平方范数为六。其最小壳同样包含
52,416,000 个向量，并为 11-design。

三个截面向量的 Gram 矩阵为

\[
H=8I_3-2J_3,\qquad H^{-1}=(I_3+J_3)/8.
\]

设 $f(y)$ 为签名 $y$ 的壳向量数。非例外域可明确写为
\begin{equation}
\begin{split}
D_0=\{y\in\mathbb Z^3:\ &|y_i|\le3,\quad |y_1+y_2+y_3|\le3,\\
&\textstyle\sum_i y_i^2+(\sum_i y_i)^2\le48\}.
\end{split}
\label{eq:domain45}
\end{equation}
截面根是 $\pm t_1,\pm t_2,\pm t_3$ 与 $\pm(t_1+t_2+t_3)$。
与这八个向量的内积界给出式~\eqref{eq:domain45}中的线性约束，
$y^{\mathsf T}H^{-1}y\le6$ 给出二次约束。它们自身的签名构成
$E=\{\pm He_1,\pm He_2,\pm He_3,\pm H(1,1,1)^{\mathsf T}\}$。
每个例外的重数为一，因为其截面投影已具有平方范数六。
枚举上述三坐标域，得到 $|D_0|=231$，故 $|D_0\cup E|=239$。

对 $D_0$ 中每对 $\{y,-y\}$ 取字典序较小的代表，令
$v_{[y]}=f(y)=f(-y)$，零向量单列。因此共有 $(231+1)/2=116$ 个变量。
对总次数为偶数且至多十的多重指标 $a$，矩阵行及右端为
\[
A_{a,[y]}=(2-\mathbf1_{y=0})\prod_i y_i^{a_i},\qquad
b_a=M_a-\sum_{y\in E}\prod_i y_i^{a_i},
\]
其中 $M_a$ 由前述球面矩公式给出。
这些单项式共有 $\sum_{j=0}^5\binom{2j+2}{2}=161$ 个，
因而完整定义了 $161\times116$ 的系统 $Av=b$。

纤维 $f(y)$ 计数满足 $(x\cdot t_1,x\cdot t_2,x\cdot t_3)=y$ 的壳向量。
目标是三个纤维的和 $f(e_1)+f(e_2)+f(e_3)$，而非确定每一个 $f(y)$。
将已知例外项移到右侧后，写矩方程为 $Av=b$，其中 $v\ge0$。
对统计量 $c^{\mathsf T}v$，若有理向量 $w$ 满足 $c-A^{\mathsf T}w\ge0$，则
\[
c^{\mathsf T}v=w^{\mathsf T}b+(c-A^{\mathsf T}w)^{\mathsf T}v
\ge w^{\mathsf T}b.
\]
这是所用矩证书的基本形式。适当选择统计量，可以让一个容易展示的纤维在下界中
带正系数，从而将寻找少量见证转化为有效的构造途径。

\begin{proposition}[纤维的正贡献]
\label{prop:moment45}
签名重数满足

\[
f(e_1)+f(e_2)+f(e_3)
\ge7377408+9f(-2,-2,3).
\]
\end{proposition}
\begin{proof}
令 $c$ 在三个类 $[e_i]$ 上的系数为 $+1$，在 $[(-2,-2,3)]$ 上为 $-9$，
其余为零。代表与多重指标均按字典序排列。
附录~\ref{app:data}所列的 \texttt{dual.json}给出161个有理乘子；
精确矩阵乘法验证 $c-A^{\mathsf T}w\ge0$ 及 $w^{\mathsf T}b=7377408$。
将其代入前述恒等式，利用 $v\ge0$ 即得结论。
\end{proof}

\begin{corollary}[45维构造]
\label{cor:section45}
指定的二次剩余偶邻格及截面给出 $K(45)\ge7380090$。
\end{corollary}
\begin{proof}
证书给出 298 个不同的整数格向量，平方范数均为六、签名均为
\((-2,-2,3)\)，因此右侧至少为 7,380,090。

将三个纤维 \(e_1,e_2,e_3\) 投影到截面的正交补，投影公式为 $\pi(x)=x-\operatorname{proj}_{U}x$，因此
\[
\pi(x)\cdot\pi(x')=x\cdot x'-y^{\mathsf T}H^{-1}y'.
\]
对 $y=e_i$，截面投影平方范数为 $1/4$；对不同 $e_i,e_j$，截面投影内积为
$1/8$。母壳中不同点的内积至多为三，故每个投影的平方范数为
\(23/4\)；同一纤维内的内积至多为 \(11/4\)，不同纤维之间至多为
\(23/8\)。除以 \(\sqrt{23/4}\) 后得到接吻构型：

\[
K(45)\ge7377408+9(298)=7380090.
\]
\end{proof}

见证搜索的一条路线在 Pless 格中得到292个向量，对应下界7380036。
扩大锚点的选择范围后，二次剩余偶邻格给出298个向量，
每增加一个见证便使投影下界增加九点。为组织这一搜索，
固定一个平方范数为八的锚点 $p\in\Lambda$，并定义
\[
A_p=\{x\in X:x\cdot p=4\}.
\]
这个母集合恰有2256点，且其完整性可以由低次矩证明。
因为 $x\pm p$ 是非零格向量，最小范数条件给出 $|x\cdot p|\le4$。
考虑多项式
\[
P(t)=t^2(t^2-1)(t^2-4)(t^2-9).
\]
它在所有允许整数值中仅于 $t=\pm4$ 非零，且 $P(4)=20160$。
八次球面矩给出 $\sum_{x\in X}P(x\cdot p)=90961920$；反极对称使
两个端点纤维等大，因此 $|A_p|=90961920/(2\cdot20160)=2256$。
又因为 $x\mapsto p-x$ 保持 $A_p$，母图的顶点是其中以 \(p/2\) 为中心的
1,128 对
\(\{x,p-x\}\)。两个点对的代表元内积为一或三时相邻；换用配对中的另一端点仍保留这个条件。重复的纤维搜索由此转为同一个母图中的共同邻点计算。

寻找锚点时，采用附录中的整数分子坐标 $p=y/\sqrt{12}$。
最终选择的 $y$ 有六个绝对值为三的坐标和42个绝对值为一的坐标，
所以 $y^2=96$。这种奇坐标类型使搜索进入邻格的另一部分，
而非只重复先前整数部分的锚点。所得到的母图度数为368，
并有一对不相邻顶点恰有70个共同邻点。差值 \(368-70=298\)
给出目标纤维。这个对应可以直接写出：将所选端点记为 $r,s$，取截面
$T=(-s,s-p,p-r)$。对与 $r$ 所在顶点相邻、与 $s$ 所在顶点不相邻的点对，
选择满足 $r\cdot a=3$ 的端点 $a$，此时 $s\cdot a=2$。
令 $w=p-a$，则
\[
\|w\|^2=6,\qquad (-s\cdot w,(s-p)\cdot w,(p-r)\cdot w)=(-2,-2,3).
\]
不同点对给出不同见证，所以298个图顶点确实对应298个所需格向量。附录\ref{app:data}中的母图检查程序重建截面，并恢复与最终证书完全一致的
298 个向量。有理矩不等式随后给出所述下界。
